\documentclass[11pt]{amsart}

%------------------------------------------------------------------------- 

% Aggiunge i bookmark ed i collegamenti nel file pdf
\usepackage[pdfauthor={Ziyang GAO}, 
        pdftitle={Enlarged mixed Shimura varieties}%
      dvips,colorlinks=true]{hyperref}
\hypersetup{
  bookmarksnumbered=true,
  linkcolor=black,
  citecolor=black,
  pagecolor=black, 
  urlcolor=black,  
}

%------------------------------------------------------------------------- 


\usepackage[utf8]{inputenc}

\usepackage{xtab}
%===Pacchetti per font matematici particolari, caratteri gotici, simboli...
\usepackage{amsmath, amssymb, cancel,verbatim}
\usepackage[all]{xy}
\usepackage[scr=euler,scrscaled=.96]{mathalfa}
 
\usepackage{enumerate}
\usepackage{mathtools,booktabs}
\usepackage{amscd}
\usepackage{bbm, stmaryrd}
\usepackage{yfonts}
\usepackage{color}

% set dimensions of columns, gap between columns, and paragraph indent 

\setlength{\textheight}{23cm}
\setlength{\textwidth}{16cm}
\setlength{\oddsidemargin}{0cm}
\setlength{\evensidemargin}{0cm}
\setlength{\topmargin}{-1cm}
%
  %====Comandi per l'inserimento di immagini
%\usepackage[pdftex]{graphicx}
\usepackage{epstopdf} 
\usepackage{booktabs}


% Con questo comando si dice a LaTex dove sono memorizzate le figure
%\graphicspath{{./Imgs/}}

% === Definizione degli ambienti theorem =================

% The Theorem environments
\setcounter{part}{-1}
\setcounter{tocdepth}{1}

%\swapnumbers
\newtheorem{teo}{Theorem}[]
\newtheorem{thm}[teo]{Theorem}
\newtheorem{ex}[teo]{Exercise}
\newtheorem{assu}[teo]{Assumption}
\newtheorem{eg}[teo]{Example}
\newtheorem{defn}{Definition}
\newtheorem{ques}[teo]{Question}
\newtheorem{rmk}[teo]{Remark}
\newtheorem{fact}{Fact}
\newtheorem{claim}{Claim}
\newtheorem{notation}[teo]{Notation}

\newcommand\numberthis{\addtocounter{equation}{1}\tag{\theequation}}
\numberwithin{equation}{section}

%
%\usepackage[small,nohug,heads=LaTeX]{diagrams}
%\diagramstyle[labelstyle=\scriptstyle]

%%%%%%%%%%% Abbreviation

%mathbb
  \newcommand{\A}{\mathbb{A}}
  \newcommand{\B}{\mathbb{B}}
  \newcommand{\C}{\mathbb{C}}
  \newcommand{\D}{\mathbb{D}}
   \newcommand{\E}{\mathbb{E}}
  \newcommand{\F}{\mathbb{F}}  
  \newcommand{\G}{\mathbb{G}}
  \renewcommand{\H}{\mathbb{H}}  
  \renewcommand{\L}{\mathbb{L}}
    \newcommand{\M}{\mathbb{M}}
  \newcommand{\N}{\mathbb{N}}
  \renewcommand{\P}{\mathbb{P}}
  \newcommand{\Q}{\mathbb{Q}}
  \newcommand{\R}{\mathbb{R}}  
  \renewcommand{\S}{\mathbb{S}}
  \newcommand{\V}{\mathbb{V}}
  \newcommand{\Z}{\mathbb{Z}}
  
  
  \renewcommand{\d}{\textbf{d}}
  \newcommand{\bH}{\textbf{H}}
  \newcommand{\bT}{\textbf{T}}
  \newcommand{\bV}{\mathbf{V}}
  
%Categories
  \newcommand{\Cop}{\mathcal{C}^{\text{op}}}
  \newcommand{\Set}{\mathbf{Set}}
  \newcommand{\SSet}{\mathbf{SSet}}
  \newcommand{\Rng}{\mathbf{Rng}}
  \newcommand{\AAlg}{\mathbf{Alg}_A}
  \newcommand{\kAlg}{\mathbf{Alg}_k}
  \newcommand{\AMod}{\mathbf{Mod}_A}
  \newcommand{\Alg}{\mathbf{Alg}}
  \newcommand{\Top}{\mathbf{Top}}
  
%Greek fonts  
  \renewcommand{\epsilon}{\varepsilon}
  
%Gothic fonts
  \newcommand{\ida}{\mathfrak{a}}
  \newcommand{\idb}{\mathfrak{b}}
  \newcommand{\idc}{\mathfrak{c}}
  \newcommand{\mi}{\mathfrak{m}} %maximal ideal m
  \newcommand{\idp}{\mathfrak{p}} %prime ideal p
  \newcommand{\idP}{\mathfrak{P}}
  \newcommand{\idq}{\mathfrak{q}} %prime ideal q
  \newcommand{\prq}{\textswab{q}} %primary ideal - q
  \newcommand{\im}{\text{Im}}
  \newcommand{\re}{\text{Re}}
  \newcommand{\tensor}{\otimes}
  \renewcommand{\cong}{\simeq}
  \renewcommand{\bar}{\overline}
  \renewcommand{\tilde}{\widetilde}
  \renewcommand{\hat}{\widehat}

  \newcommand{\epi}{\varepsilon}
  \providecommand{\loc}[1]{S^{-1}#1}
  \providecommand{\frac}[1]{\mathrm{Frac}(#1)}
  \newcommand{\spec}{\mathrm{Spec}}
  \providecommand{\nil}[1]{\mathrm{Nil}(#1)}
  \providecommand{\SqMat}[2]{{\mathrm{M}}_{#2, #2}(#1)}
  \renewcommand{\hom}{\mathrm{Hom}}
  \newcommand{\tor}{\mathrm{Tor}}
  \newcommand{\ext}{\mathrm{Ext}}
  \newcommand{\proj}{\mathrm{proj}}
  \newcommand{\inj}{\mathrm{inj}}
  \newcommand{\depth}{\mathrm{depth}}
  \newcommand{\supp}{\mathrm{Supp}}
  \renewcommand{\proj}{\mathrm{Proj}}
  \newcommand{\sym}{\mathrm{Sym}}
  \newcommand{\Spec}{\mathrm{\textbf{Spec}}}
  \newcommand{\Proj}{\mathrm{\textbf{Proj}}}
  \newcommand{\Frac}{\mathrm{Frac}}
  \newcommand{\rank}{\mathrm{rank}}
  \newcommand{\cl}{\mathrm{Cl}}
  \newcommand{\gal}{\mathrm{Gal}}
  \newcommand{\reg}{\mathrm{Reg}}
  \newcommand{\frob}{\mathrm{Frob}}
  \newcommand{\stab}{\mathrm{Stab}}
  \newcommand{\MT}{\mathrm{MT}}
  \newcommand{\GL}{\mathrm{GL}}
  \newcommand{\SL}{\mathrm{SL}}
  \newcommand{\Sh}{\mathrm{Sh}}


    \newcommand{\vol}{\mathrm{Vol}}  
  \newcommand{\tr}{\mathrm{Tr}}  
  \newcommand{\ann}{\mathrm{Ann}}
  \newcommand{\coker}{\mathrm{Coker}}
  \renewcommand{\ker}{\mathrm{Ker}}
  \newcommand{\ch}{\mathrm{Ch}}
  \newcommand{\Max}{\mathrm{Max}}
    \newcommand{\height}{\mathrm{ht}} %height of a prime
  \newcommand{\coht}{\mathrm{coht}}  %depth of a prime
  \newcommand{\codim}{\mathrm{codim}}
  \newcommand{\Mat}{\mathrm{M}}
  \newcommand{\Trdeg}{\mathrm{Trdeg}}
  \newcommand{\colim}{\mathrm{colim}}
  \renewcommand{\lim}{\mathrm{lim}}
  \renewcommand{\deg}{\mathrm{deg}}
  \renewcommand{\div}{\mathrm{div}}
  \newcommand{\ord}{\mathrm{ord}}
  \newcommand{\pic}{\mathrm{Pic}}
  \newcommand{\aut}{\mathrm{Aut}}
  \newcommand{\End}{\mathrm{End}}
  \newcommand{\Int}{\mathrm{Int}}
  \newcommand{\disc}{\mathrm{disc}}
  \newcommand{\lie}{\mathrm{Lie}}
  \newcommand{\Gr}{\mathrm{Gr}}
  \newcommand{\Alb}{\mathrm{Alb}}
  \newcommand{\alb}{\mathrm{alb}}
  \newcommand{\der}{\mathrm{der}}
  \newcommand{\ad}{\mathrm{ad}}
  \newcommand{\sm}{\mathrm{sm}}
  \newcommand{\de}{\mathrm{de}}
  \newcommand{\n}{\mathrm{n}}
  \newcommand{\cp}{\mathrm{cp}}
  \newcommand{\mon}{\mathrm{mon}}
  \newcommand{\Zar}{\mathrm{Zar}}
  \newcommand{\GSp}{\mathrm{GSp}}
  \newcommand{\Sp}{\mathrm{Sp}}
  \newcommand{\AL}{\mathrm{AL}}
  \renewcommand{\Re}{\mathfrak{R}\mathfrak{F}}
  \newcommand{\ab}{\mathrm{ab}}
  \newcommand{\res}{\mathrm{Res}}
  \newcommand{\rec}{\mathrm{rec}}
  \newcommand{\CM}{\mathrm{CM}}
  \newcommand{\unif}{\mathrm{unif}}
  \newcommand{\Pic}{\mathrm{Pic}}
  \newcommand{\Div}{\mathrm{Div}}

%Other
\newcommand{\cA}{\mathcal{A}}
\newcommand{\cB}{\mathcal{B}}
\newcommand{\cC}{\mathcal{C}}
\newcommand{\cD}{\mathcal{D}}
\newcommand{\cE}{\mathcal{E}}
\newcommand{\cF}{\mathcal{F}}
\newcommand{\cG}{\mathcal{G}}
\newcommand{\cH}{\mathcal{H}}
\newcommand{\cI}{\mathcal{I}}
\newcommand{\cJ}{\mathcal{J}}
\newcommand{\cK}{\mathcal{K}}
\newcommand{\cL}{\mathcal{L}}
\newcommand{\cM}{\mathcal{M}}
\newcommand{\cO}{\mathcal{O}}
\newcommand{\cR}{\mathcal{R}}
\newcommand{\cS}{\mathcal{S}}
\newcommand{\cT}{\mathcal{T}}
\newcommand{\cU}{\mathcal{U}}
\newcommand{\cV}{\mathcal{V}}
\newcommand{\cW}{\mathcal{W}}
\newcommand{\cX}{\mathcal{X}}
\newcommand{\cY}{\mathcal{Y}}
\newcommand{\cZ}{\mathcal{Z}}

\newcommand{\bpi}{\boldsymbol{\pi}}
\newcommand{\brho}{\boldsymbol{\rho}}


\renewcommand*{\thefootnote}{[\arabic{footnote}]}


%dotted curved arrow above
\newcommand{\upperarrow}[1]{
\setlength{\unitlength}{0.03\DiagramCellWidth}
\begin{picture}(0,0)(0,0)
\linethickness{1pt} 
\bezier{15}(-46,9)(-24,23)(-2,9)
\put(-24,19){\makebox(0,0)[b]{$\scriptstyle {#1}$}}
\put(-46.6,8.7){\vector(-2,-1){0}}
\end{picture}
}
\newcommand{\upperarrowtwo}[1]{
\setlength{\unitlength}{0.03\DiagramCellWidth}
\begin{picture}(0,0)(0,0)
\linethickness{1pt} 
\bezier{15}(-40,9)(-21,23)(-2,9)
\put(-21,19){\makebox(0,0)[b]{$\scriptstyle {#1}$}}
\put(-40.6,8.7){\vector(-2,-1){0}}
\end{picture}
}



%dotted curved arrow on the left
\newcommand{\leqslantftcurvedarrow}[1]{
\setlength{\unitlength}{0.03\DiagramCellWidth}
\begin{picture}(0,0)(0,0)
\linethickness{1pt} 
\bezier{15}(-9,-16)(-23,5)(-9,26)
\put(-27,5){\makebox(0,0)[b]{$\scriptstyle {#1}$}}
\put(-7.6,26){\vector(2,1){0}}
\end{picture}
}

\newcommand\supervisor[1]{\def\@supervisor{#1}}


\newcommand{\pullback}[1][dr]{\save*!/#1-1.2pc/#1:(-1,1)@^{|-}\restore}
\newcommand{\bigpullback}[1][dr]{\save*!/#1-1.7pc/#1:(-1.5,1.5)@^{|-}\restore}
\newcommand{\pushoutcorner}[1][dr]{\save*!/#1+1.7pc/#1:(1.5,-1.5)@^{|-}\restore}


\newcounter{elno}   
\newenvironment{romanlist}{
                         \begin{list}{\roman{elno})
                                     }{\usecounter{elno}}
                      }{
                         \end{list}}
                         
                         
\newcommand{\mail}[1]{\href{mailto:#1}{\texttt{#1}}}

%===MACRO FOR ARROWS ====
% Surjective Arrow
\newcommand{\surj}{\twoheadrightarrow}
\newcommand{\into}{\hookrightarrow}
\renewcommand{\cong}{\simeq}

\newcommand{\longby}[1]{\stackrel{#1}{\longrightarrow}}

\begin{document}

\section*{Conference Program}

\subsection*{Monday, August 3}

\subsubsection*{9:30--10:30 \quad Xinyi Yuan}

\textbf{On Vojta's Proof of the Mordell Conjecture}

The Mordell conjecture was originally proved by Faltings in 1983, and a second proof was given by Vojta in 1989. In this talk, I will describe Vojta's proof with a simplification by using my arithmetic Siu inequality.

\bigskip
\hrule
\bigskip

\subsubsection*{11:00--12:00 \quad Isabel Vogt}

\textbf{Unlikely Ramification of Algebraic Points on Curves}

A slogan of arithmetic geometry is that ``geometry controls arithmetic'': as the geometric complexity increases, the arithmetic also becomes more complicated. In this talk, I will discuss results in this direction that show that there are many number fields that cannot appear as the residue field of points on a fixed curve of genus at least 2. This is joint work with Bianca Viray.

\bigskip
\hrule
\bigskip

\subsubsection*{14:00--15:00 \quad Jackson Morrow}

\textbf{A non-Archimedean Ax--Lindemann theorem for certain $p$-adically uniformized Shimura varieties}

The Ax--Lindemann conjecture is a functional transcendence statement in arithmetic geometry concerning the uniformization map of an arithmetic variety. Klingler--Ullmo--Yafaev proved this conjecture for the complex analytification of a Shimura variety. I will present a non-Archimedean analogue of this result for certain $p$-adically uniformized Shimura varieties, in particular quaternonic Shimura varieties and a class of PEL-type unitary Shimura varieties treated by Rapoport--Zink. Our proof has two main ingredients. First, we prove that the complex and $p$-adic uniformization maps for these Shimura varieties satisfy the same system of differential equations --- this result is an extension of a theorem of Andr'e from the setting of Shimura curves to our broader setting. Second, we use techniques from differential algebra and results in model theory concerning the theory of differential closed fields to show that the validity of complex analytic Ax--Lindemann statement is equivalent to the validity of $p$-adic analytic Ax--Lindemann statement in our setting. This is joint work with Marc-Hubert Nicole and Giovanni Rosso.

\bigskip
\hrule
\bigskip

\subsubsection*{15:30--16:30 \quad Alex Betts}

\emph{Title and abstract to be announced.}

\newpage

\subsection*{Tuesday, August 4}

\subsubsection*{9:30--10:30 \quad Jit Wu Yap}

\textbf{Uniform boundedness of torsion points for abelian varieties over function fields}

Let $K$ be the function field of a curve $B$ over $\mathbb{C}$ and $A/K$ an abelian variety with trivial trace. The uniform boundedness conjecture predicts that the number of torsion points in $A(K)$ is bounded solely in terms of $\dim A$ and $K$. In this talk, I will present a proof of this conjecture and also discuss some related questions. This is joint work with Nicole Looper.

\bigskip
\hrule
\bigskip

\subsubsection*{11:00--12:00 \quad Riccardo Pengo}

\textbf{Explicit bounds and algorithms to find torsion points}

A celebrated conjecture by Manin and Mumford, now a theorem, predicted the distribution of torsion points in subvarieties of (abelian) group varieties. Our talk will be focused on the explicit aspects surrounding these conjectures, with the objective of finding the aforementioned torsion points in explicit infinite families of varieties. This talk will be based on joint work in progress with Evelina Viada.

\bigskip
\hrule
\bigskip

\subsubsection*{14:00--15:00 \quad Boya Wen}

\textbf{Local systems over graphs and CM cycles over Shimura Curves}

A Shimura Curve parametrizes principally polarized abelian surfaces with quaternion multiplication by a quaternion algebra $B$. At a prime where $B$ ramifies, the Shimura curve reduces to a split degenerate curve whose dual graph is a finite regular graph of degree $p+1$. In this talk, we will introduce local systems of spheres and inner product spaces on such graphs, as well as harmonic analysis on them. We will also connect these constructions to arithmetic intersections of CM cycles in a universal abelian surface above the Shimura curve.

\bigskip
\hrule
\bigskip

\subsubsection*{15:30--16:15 \quad Lightning Talks}
(Hanson Hao, Zhelun Chen, Xinyu Fang, Dan 	Townsend, Xiao Zhong, Yuta Nakayama, Soheil Memarian)

\bigskip
\hrule
\bigskip

\subsubsection*{16:30--17:30 \quad Shouwu Zhang}

\textbf{Adelic Line Bundles: From Canonical Heights to Arithmetic Positivity}

Canonical heights play a fundamental role in arithmetic geometry, from the Mordell--Weil theorem to the Birch--Swinnerton-Dyer conjecture. In 1995, I introduced the notion of adelic line bundles, providing a geometric interpretation of canonical heights as degrees of metrized line bundles over all places of a number field. This viewpoint unifies Archimedean and non-Archimedean geometry and brings analytic methods into arithmetic geometry.

In this lecture, I will explain how adelic line bundles lead naturally to adelic intersection theory, Monge--Amp\`ere measures, and the equidistribution of small points, with applications to the Bogomolov conjecture and arithmetic dynamics. I will then describe recent joint work with Xinyi Yuan that extends the theory to quasi-projective varieties, thereby making it possible to construct canonical adelic metrics on moduli spaces. These developments reveal new positivity phenomena in arithmetic geometry and provide a conceptual framework for studying heights of algebraic cycles and related arithmetic questions.

\newpage

\subsection*{Wednesday, August 5}

\subsubsection*{9:00--10:00 \quad Myrto Mavraki}

\textbf{Special Subvarieties in Families of endomorphisms of $(\mathbb{P}^1)^n$}

Let $\Phi=(f_1,\ldots,f_n)$ be an algebraic family of coordinatewise rational maps on $(\mathbb P^1)^n$, and let $X$ be a family of subvarieties. When is $X$ dynamically special---for instance, preperiodic---for a Zariski-dense set of parameters? Inspired by Pink's conjectures in arithmetic geometry and recent work of Gao--Habegger, Yuan and Zhang, I will discuss a geometric answer to this question. I will explain its connections with bifurcation theory in complex dynamics and the positivity of canonical adelic line bundles, and indicate applications to uniformity questions for preperiodic points and invariant subvarieties. This is joint work with Laura DeMarco.

\bigskip
\hrule
\bigskip

\subsubsection*{10:30--11:30 \quad Junyi Xie}

\textbf{Abelian points in backward orbits}

Let $K$ be a number field and let $f$ be a polarized self-map of a normal projective variety $X$. For a point $x$ in $X$, consider the field generated by all iterated preimages of $x$ under $f$. I will explain when this field extension is virtually abelian. The answer is exactly the exceptional case: $f$ must be an exceptional map with complex multiplication, and $x$ must be preperiodic. This proves a conjecture of Andrews and Petsche and extends it to higher dimensions.

\bigskip
\hrule
\bigskip

\subsubsection*{12:00--13:00 \quad Gabriel Dill}

\textbf{Intersecting a curve with multiples of another curve}

I will present joint work in progress with Fabrizio Barroero and Lars K\"uhne on the abelian analogue of a question of Levin, which has appeared in Zannier's book on unlikely intersections more than a decade ago. Namely, we study the locus of points that lie on a given curve and have an integer multiple lying on another given curve.

\newpage

\subsection*{Thursday, August 6}

\subsubsection*{9:30--10:30 \quad Jos\'e Burgos Gil}

\textbf{The essential minimum of height functions on the projective line}

The essential minimum of a height function is the minimal value that the height function can attain at generic points. There are methods to compute upper and lower bounds for the essential minimum. In a joint work with Binggang Qu, Ricardo Menares, and Martin Sombra, we prove using linear programming techniques that the difference between the upper and lower bounds can be made arbitrarily small. Therefore, one can devise a theoretical algorithm to compute the essential minimum with arbitrary precision and thus the essential minimum is a ``computable'' real number. This result has applications in several classical problems like the integral Chebyshev constant of the unit interval, the spectrums of the Zhang-Zagier and the Faltings heights and the asymptotic behaviour of the length of the shortest vector in the lattice associated with the Grassmannian $\operatorname{Gr}(2,4)$.

\bigskip
\hrule
\bigskip

\subsubsection*{11:00--12:00 \quad Klaus K\"unnemann}

\textbf{A tropical formula for non-archimedean local heights}

We report on joint work with Jos\'e Burgos and Walter Gubler. Let $X$ be a smooth projective variety over a non-archimedean field. Using tropicalization, one can introduce real-valued forms and currents on the non-archimedean analytification of $X$. We discuss how these can be used to compute non-archimedean local heights.

\bigskip
\hrule
\bigskip

\subsubsection*{14:00--15:00 \quad Alice Lin}

\textbf{Finiteness of heights in isogeny classes of motives}

Using integral $p$-adic Hodge theory, Kato and Koshikawa define a generalization of the Faltings height of an abelian variety to motives defined over a number field. Assuming the adelic Mumford-Tate conjecture, we prove a finiteness property for heights in the isogeny class of a motive, where the isogenous motives are not required to be defined over the same number field. This expands on a result of Kisin and Mocz for the Faltings height in isogeny classes of abelian varieties.

\bigskip
\hrule
\bigskip

\subsubsection*{15:30--16:15 \quad Lightning Talks}
(Ying Wang, Wenbin Luo, Michal Szachniewicz, Zheng Xiao, Mithun Padinhare, Yu Fu)
\bigskip
\hrule
\bigskip

\subsubsection*{16:30--17:30 \quad Joseph H. Silverman}

\textbf{Positivity Properties of Fourier Expansions of Local Heights on Abelian Varieties}

Let $K$ be a complete local field, and let $A/K$ be an abelian variety of dimension $g$. If $A$ has good reduction, then the intersection pairing $P\cdot D$ on the N'eron model of $A$ gives a local canonical height $A(K)\smallsetminus|D|\to\mathbb{R}$ associated to the ample effective divisor $D$. If $A$ has split multiplicative reduction, there is a Bernoulli correction term that takes on both positive and negative values. In the case $g=1$, this term is given by the periodic $2$nd Bernoulli polynomial $\mathbb{B}_2(T)=T^2-T+\frac16$ whose classical Fourier expansion has non-negative coefficients. For $g=2$, Nicole Looper and I investigated the analogous $2$-variable periodic Bernoulli polynomial $\mathbb{B}_2(X,Y)$. We gave explicit formulas for the Fourier expansion of $\mathbb{B}_2(X,Y)$ and showed how weighted averages could be used to eliminate the negative Fourier coefficients. In this talk I will describe the $g=1$ and $g=2$ cases, briefly mention some motivating applications, and raise some questions in higher dimensions.

\bigskip
\hrule
\bigskip

\subsubsection*{17:40 \quad Reception}

\newpage

\subsection*{Friday, August 7}

\subsubsection*{9:00--10:00 \quad Philipp Habegger}

\textbf{Specializing Linear Recurrence Sequences at Roots of Unity}

The Skolem-Mahler-Lech Theorem characterizes the vanishing members of a linear recurrence sequence. Recently, Bacik gave an effective proof of this classical theorem for four term linear recurrence sequences in the number field case. For orders greater than four, no effective proof is known. In other words, given a general linear recurrence sequence, we know of no algorithm that is guaranteed to determine the indices at which the sequence vanishes.

We consider linear recurrence sequences of rational functions and study when sequence members vanish at a root of unity. Bilu-Luca and Ostafe-Shparlinski proved finiteness results for linear recurrence sequences of order two. I will report on a new finiteness result for linear recurrence sequences of order three. We require linear forms in logarithms combined with the Pila-Zannier strategy as well as the function field $abc$ Conjecture.

This is joint work in progress with Alina Ostafe and David Masser.

\bigskip
\hrule
\bigskip

\subsubsection*{10:30--11:30 \quad Yunqing Tang}

\textbf{The arithmetic of power series and applications to irrationality}

We will first briefly discuss our approach to prove irrationality of certain periods such as certain product of two log values. The key ingredient is an arithmetic holonomy theorem built upon earlier work by Andr\'e, Bost, Charles (and others) on arithmetic algebraization theorems via Arakelov theory. We will then discuss our result on irrationality measures and a proof of transcendence of $\pi$ in our framework. This is joint work with Frank Calegari and Vesselin Dimitrov.

\bigskip
\hrule
\bigskip

\subsubsection*{12:00--13:00 \quad Aaron Levin}


\textbf{Some applications of Diophantine approximation to integral points and unlikely intersections}

We discuss recent results in Diophantine approximation, with a focus on heights attached to closed subschemes of higher codimension and greatest common divisor type inequalities. We then give some applications to integral points on surfaces, and discuss relations between GCD inequalities and a new unlikely intersection conjecture. This is joint work with Zheng Xiao and Keping Huang.


\end{document}